\chapter{How to Use This Book}

\section*{One idea}
Nobody can study every endgame. Endgame manuals are thick, and most of what they teach you will
rarely need. This book starts from the other end. It asks where players of one rating band
actually lose points in the endgame, and studies those endgames first. The band is named on the
cover.

To find out, every endgame in thousands of over-the-board tournament games of that band was
checked by a chess engine. Each time a player turned a won position
into a draw, or a drawn position into a loss, the move was marked as a mistake. Counting those
mistakes shows which endgames are worth your time.

\section*{\xref{Part}{part:mistakes}: where the mistakes are}
\xref{Chapter}{ch:meet} shows which endgames you will meet most often. \xref{Chapter}{ch:go-wrong} shows which ones go wrong
most often once they are on the board. \xref{Chapter}{ch:where-to-start} puts the two together: the endgames at the top
of its list cost the most points, and they set the order of the rest of the book.

\section*{\xref{Part}{part:endgames}: the endgames}
Each chapter in \xref{Part}{part:endgames} takes one endgame from that list. It shows positions in which a player
went wrong, from the simplest to the more complex, each followed by the rest of the game.

Use them as exercises. Cover the moves below a diagram and decide what you would play. Then
read on and compare. Every diagram has a link that opens the position on Lichess, where you can
play it out against the computer or check your ideas.

\section*{Why mistakes from your own band?}
The games come from players of one rating band, so the mistakes are the ones players at that
level really make. If you understand why they went wrong, you will avoid the same mistakes in
your own games.
